% ===========================================================================
%  sviluppatore/05-estendere/08-fonte.tex
% ===========================================================================
\section{Una fonte di dati: tenere aggiornata una categoria}\label{sez:nuova-fonte}

Una \emph{categoria} è una cartella di \pkgDati{} (Settings › Data › Data
categories): \cod{earnings}, \cod{fundamentals}, \cod{news}, \cod{rates},
\cod{macro}, \cod{mms}, o una creata dall'utente. Un componente dichiara
quali legge in \cod{required\_external\_data} e ogni quanto vanno aggiornate
in \cod{maintenance}. La \emph{fonte} è il codice che sa procurare i dati di
una categoria e li tiene aggiornati: \file{data/fonti/} per quelle di
sistema, il file \file{data/<categoria>/\_source.py} per le altre.

\subsection{Il buco che chiude}

Fino a questa versione le dichiarazioni c'erano e nessuno le onorava. Il
catalogo di \file{data/external\_data.py} scaricava solo ciò che mancava del
tutto --- le notizie solo se non ce n'era nessuna, i fondamentali dei simboli
mai visti, i tassi mai --- e \cod{earnings} e \cod{macro} non li conosceva.
Le porte sui dati delle notizie (\file{data\_port.py} di
\cod{llm\_news\_macro} e \cod{news\_pressure}) chiamavano una funzione che,
trovato un archivio qualsiasi, non faceva nulla; ora non ci sono più, perché
una porta che scrive nella cartella di una categoria è un secondo scrittore
con regole sue. Dal vivo i file
restavano fermi al giorno del primo download: una strategia sugli utili non
vedeva il trimestre nuovo, e nessun messaggio lo diceva.

\subsection{Le fonti di sistema}

\begin{center}
\begin{tabular}{@{}L{26mm}L{60mm}L{36mm}L{14mm}@{}}
\toprule
\textbf{Categoria} & \textbf{Fonte} & \textbf{Chiave} & \textbf{Ogni} \\
\midrule
\rigaTenue
\cod{earnings} & Polygon, bilanci trimestrali (EPS). Chiede solo quando un
  trimestre nuovo è atteso: \num{70} giorni dopo l'ultimo deposito, poi una
  volta al giorno & \cod{POLYGON\_API\_KEY} & 1 g \\
\cod{fundamentals} & Polygon, stesso endpoint e stessa regola (la risposta
  si riusa nello stesso giro) & \cod{POLYGON\_API\_KEY} & 90 g \\
\rigaTenue
\cod{news} & Polygon; GDELT senza chiave. Riparte dal giorno dell'ultima
  notizia del file & --- & 1 g \\
\cod{rates} & FRED, tasso in vigore per valuta (otto valute) &
  \cod{FRED\_API\_KEY} & 7 g \\
\rigaTenue
\cod{macro} & FRED (curva, spread, lavoro, CPI) e Yahoo (VIX, solo a seduta
  chiusa); estende anche il risk-off macro & \cod{FRED\_API\_KEY} & 1 g \\
\cod{mms} & calcolate dall'app sulle barre chiuse (heatmap di liquidità) &
  --- & 1 g \\
\bottomrule
\end{tabular}
\end{center}

\noindent
Le chiavi si leggono dall'\file{.env} dell'app (\simbolo{dashboard.envload}).
La chiave Polygon della pagina Data resta nel browser e non serve agli
aggiornamenti automatici. Un \file{\_source.py} in una categoria di sistema
prende il posto della fonte di sistema: gli utili da un altro provider si
scrivono lì.

\subsection{Il contratto}

\begin{codiceInterno}{python}{La fonte di una categoria dell'utente}{data/sentiment/\_source.py}
import os
import pandas as pd
import requests

DESCRIPTION = "Acme Data - daily sentiment"   # pagina Data e registro del runner
REQUIRES = ["ACME_API_KEY"]                   # nell'.env dell'app
EVERY_DAYS = 1                                # se la strategia non dichiara la cadenza
KEY = "index"                                 # o una lista di colonne
AFTER_ONLY = True                             # serie: solo dopo l'ultima data

def fetch(name, since, until, *, base_path, log):
    start = "2015-01-01"
    if since is not None:
        start = (since + pd.Timedelta(days=1)).date().isoformat()
    r = requests.get("https://api.acme.example/v1/sentiment",
                     params={"symbol": name, "from": start, "to": until,
                             "key": os.environ["ACME_API_KEY"]}, timeout=30)
    r.raise_for_status()
    df = pd.DataFrame(r.json()["data"])
    if df.empty:
        return df
    df["date"] = pd.to_datetime(df["date"])
    return df.set_index("date")[["score"]].sort_index()
\end{codiceInterno}

\noindent
\cod{fetch} riceve il nome del file senza estensione (un simbolo nelle
categorie 1:1, un file condiviso nelle 1:N), \cod{since} --- l'ultima data
che il file ha già, \cod{None} se il file non c'è --- e \cod{until}, la data
fin dove servono i dati. Restituisce righe: anche righe che il file ha già,
che non entrano due volte. Tre funzioni facoltative:

\begin{center}
\begin{tabular}{@{}L{26mm}L{112mm}@{}}
\toprule
\textbf{Funzione} & \textbf{A che serve} \\
\midrule
\rigaTenue
\cod{files} & \cod{(base\_path, symbols, context)}: quali file aggiornare.
  Senza: i simboli dei componenti nelle 1:1 (tutti i file della cartella da
  \voce{Update now}), i file presenti nelle 1:N. \cod{context} è
  \cod{live}, \cod{backtest} o \cod{manual}. \\
\cod{due} & \cod{(name, since, info, now)}: se vale la pena chiedere adesso;
  gli utili la usano per non chiedere ogni giorno quattrocento bilanci che
  non sono cambiati. \cod{info["controllato"]} è l'ultimo controllo del
  file. \\
\rigaTenue
\cod{after\_update} & \cod{(base\_path, result, log)}: ciò che dipende
  dalle righe nuove; la fonte \cod{macro} estende il risk-off macro. \\
\bottomrule
\end{tabular}
\end{center}

\subsection{La regola dei file: si aggiunge, non si riscrive}

Una fonte non scrive mai un file da sé. Le righe passano da
\simbolo{data.estendi.estendi}, che aggiunge soltanto quelle la cui
\emph{chiave} il file non ha ancora:

\begin{description}[leftmargin = 1.4em, style = nextline, itemsep = 4pt]
  \item[\cod{KEY = "index"}] serie datate, indice di date. Con
    \cod{AFTER\_ONLY} entrano solo le date dopo l'ultima del file.
  \item[\cod{KEY = ["filing\_date", "period\_end"]}] eventi: una riga entra
    se il file non ne ha una con gli stessi valori in quelle colonne, anche
    con una data passata --- un deposito che il provider pubblica in ritardo
    era pubblico in quel giorno --- e va al suo posto di data.
\end{description}

\begin{perche}{Perché non «riscarico tutto e sovrascrivo»}
Tre ragioni. Un backtest rifatto domani vede, per i giorni passati, gli
stessi numeri di oggi. Un dato corretto a posteriori dal provider (un utile
riesposto, una serie rivista) non riscrive ciò che si sapeva quel giorno, che
è il modo onesto di tenere un archivio che serve anche dal vivo. Un download
andato storto a metà non cancella nulla: al peggio aggiunge meno righe.
\end{perche}

\noindent
Il file tiene il suo formato: un CSV si legge come testo e si riscrive con la
stessa intestazione, lo stesso separatore e lo stesso stile delle date; un
parquet con il suo indice. La scrittura è atomica (file temporaneo e
\cod{os.replace}) sotto un blocco per file in \file{.lock/}, che funziona
anche fra processi: due runner, o un runner e il pulsante, si mettono in fila
invece di perdere righe a vicenda. Le serie che FRED pubblica una volta al
mese si salvano alla data dell'osservazione e solo fino all'ultima: il mese
che esce in ritardo entra con le sue date al primo aggiornamento utile. Il
ritardo di pubblicazione lo applica chi legge (\cod{macro\_regime}).

\subsection{Quando gira}

\begin{center}
\begin{tabular}{@{}L{26mm}L{112mm}@{}}
\toprule
\textbf{Dove} & \textbf{Che cosa succede} \\
\midrule
\rigaTenue
live & il runner calcola il piano all'avvio e lo scrive nel registro; a ogni
  giro avvia in sottofondo le categorie scadute, senza fermare il ciclo
  (\vediSez{live-dati}). Le righe nuove vanno in
  \file{data/live/<categoria>/}, non nei file dei backtest
  (\vediSez{live-dati-vivo}). \\
backtest & all'avvio, \emph{prima} di \cod{prepare\_extra\_data}: ogni
  categoria dichiarata si estende fino alla fine del periodo; se un altro
  lavoro la sta aggiornando, si aspetta che finisca. \\
\rigaTenue
pagina Data & \voce{Update now}, un job su tutti i file della categoria. \\
\bottomrule
\end{tabular}
\end{center}

\noindent
Quando e come è andato ogni aggiornamento sta in
\file{data/.aggiornamenti.json}: comune a tutti i runner e ai riavvii, letto
dalla pagina Data. Mai fatto: scaduto. Andato bene: alla cadenza. Andato male
o chiave mancante: si riprova dopo un'ora (o alla cadenza, se più breve).
\simbolo{fonti.versione} cambia ogni volta che arrivano righe nuove: è così
che il motore live sa quando rifare i dati extra di una strategia.

\begin{attenzione}{Due cose che non si restituiscono mai}
Un valore che diventa noto dopo la sua data va restituito con la data in cui
è diventato noto (la pubblicazione), non con quella a cui si riferisce:
altrimenti il backtest lo vede prima del tempo. E una barra ancora in corso
(il prezzo di oggi a seduta aperta) non si restituisce: una volta aggiunta ci
resta, come se fosse la chiusura. Per questo il VIX da Yahoo entra solo dopo
le 16:30 di New York e le feature MMS solo sulle barre chiuse.
\end{attenzione}

\begin{pratica}{Dare una fonte a una categoria}
\begin{enumerate}
  \item Pagina Data, card della categoria: \voce{Create source file}. Crea
        \file{data/<categoria>/\_source.py} dal modello
        (\file{data/fonti/modello.py}) e non sovrascrive un file che c'è.
  \item Scrivere \cod{fetch}; le chiavi nell'\file{.env} dell'app, mai nel
        file.
  \item \voce{Update now}: il job dice quante righe sono entrate e quali
        file sono falliti; la card mostra l'esito.
  \item Una fonte che non si carica (un errore di sintassi) compare in rosso
        sulla card e nel piano del runner, e si riprova solo quando il file
        cambia.
\end{enumerate}
\end{pratica}

\begin{attenzione}{\cod{until} non è sempre oggi}
L'unione dei dati del live con il deposito (\vediSez{live-dati-vivo}) cerca i
periodi vuoti di una serie --- i giorni in cui l'app era spenta --- e li
chiede alla fonte: \cod{fetch(nome, since, until)} con \cod{since} il giorno
prima del vuoto e \cod{until} il suo ultimo giorno. Una fonte che onora solo
\cod{since} restituisce fino a oggi, e va bene (le righe di troppo si
scartano); una che ignora entrambi i limiti e restituisce sempre le ultime
trenta righe non colma nulla, e l'utente si trova a decidere un buco che la
fonte avrebbe potuto riempire.
\end{attenzione}

\noindent
Il riferimento completo, in inglese, è \file{DATA\_SOURCE\_TEMPLATE.md}. La
prova \file{tests/prova\_fonti\_dati.py} copre ogni fonte con provider finti.
